\documentclass[11pt]{letter}
\usepackage[margin=1in]{geometry}
\usepackage{hyperref}

\signature{Wen-Hsien Lin, Ph.D.\\AI and Data Applications Division\\GGA Corp.\\Taipei, Taiwan\\BryceLin@bionetTX.com}

\address{Wen-Hsien Lin\\AI and Data Applications Division\\GGA Corp.\\Taipei, Taiwan}

\begin{document}

\begin{letter}{Editorial Office\\BMC Bioinformatics}

\opening{Dear Editor,}

We are pleased to submit our manuscript entitled ``\textbf{DeepExoMir: A Hybrid Deep Learning Framework Integrating RNA Language Model Embeddings and Duplex Graph Attention for MicroRNA Target Prediction}'' for consideration as a Methodology Article in \textit{BMC Bioinformatics}.

Accurate prediction of microRNA targets remains a fundamental challenge in computational RNA biology. Recent benchmarking with experimentally validated CLIP-seq negatives (miRBench) has revealed that many widely used methods perform substantially worse than originally reported, highlighting the need for methods evaluated under rigorous conditions.

Our manuscript presents DeepExoMir, which makes four key methodological contributions:

\begin{enumerate}
\item \textbf{RNA language model integration:} We demonstrate that frozen embeddings from the RiNALMo foundation model (650M parameters), combined with PCA dimensionality reduction, provide effective sequence representations without task-specific fine-tuning.

\item \textbf{Duplex graph attention network (DuplexGAT):} To our knowledge, this is the first application of graph attention networks to explicitly model the miRNA-target duplex as a nucleotide-level graph with typed edges for backbone connectivity, base-pairing, and stacking proximity.

\item \textbf{Hybrid encoder architecture:} An 8-layer encoder alternating bidirectional convolution gating with cross-attention, paired with interaction-aware pooling and mixture-of-experts classification.

\item \textbf{Comprehensive ablation:} Systematic evaluation across 22 model versions quantifies each component's contribution, revealing that evolutionary conservation features account for 52\% of total feature importance.
\end{enumerate}

DeepExoMir achieves a mean AU-PRC of 0.86 on the miRBench benchmark, surpassing the best retrained CNN baseline (0.82) by 0.04. Multi-seed evaluation confirms high reproducibility (test AUC SD = 0.001). Attention analysis reveals biologically meaningful patterns including seed region enrichment and 3$'$ compensatory pairing.

All source code, pre-trained models, and analysis scripts are publicly available at \url{https://github.com/linwenh09/DeepExoMir} and archived at Zenodo (\url{https://doi.org/10.5281/zenodo.19216306}).

This manuscript has not been published previously and is not under consideration elsewhere. The author confirms that this work is original and declares no competing interests.

We believe this work is well suited for the scope and readership of \textit{BMC Bioinformatics}, as it directly builds upon methods previously published in the journal and provides a rigorous methodological contribution to the miRNA target prediction field.

Thank you for your consideration.

\closing{Sincerely,}

\end{letter}
\end{document}
